\begin{resumo}[Abstract]
\begin{otherlanguage*}{english}
\begin{SingleSpace}
The use of \textit{smartphones} has increased over the years, becoming an essential part of daily life. Despite the benefits of connectivity, excessive use of these devices is associated with physical, mental, cognitive, and social impacts. Among the main problems reported in the literature are sleep disturbances, anxiety, attention difficulties, neck pain, and headaches. Existing solutions, such as native app blocking features, are easily bypassed by users themselves. This work proposes the development of the Desconecta application, with a gamified approach in which groups of users participate in weekly challenges to maintain lower screen time within a given period. The application allows users to monitor their screen time, set goals, and compare results among participants. The development process included user research for data collection and requirements gathering, followed by the implementation of the solution and its evaluation with the target audience. In the tests conducted, a reduction in the average usage time among participants and good engagement in the proposed dynamics were observed, indicating that the collective approach may contribute to a more conscious and healthy use of technology.
\end{SingleSpace}

%Eventually you can also write it in spanish \textit{(resumen}), french \textit{(résumé)}, italian \textit{(riassunto)} etc.

\vspace{\onelineskip}
   \textbf{Keywords}: excessive \textit{smartphone} use; screen time; mental health; persuasive \textit{design}; control application.
 \end{otherlanguage*}
\end{resumo}